\documentclass{article}
\usepackage{amsmath,amssymb}
\usepackage{xurl}
\usepackage[hidelinks]{hyperref}
\setlength{\emergencystretch}{2em}
% Shared commands for the notes in docs/paper-gaps/.
% Every note loads this file via % Shared commands for the notes in docs/paper-gaps/.
% Every note loads this file via % Shared commands for the notes in docs/paper-gaps/.
% Every note loads this file via \input{command}. Project-wide notation
% belongs here, never in per-note redefinitions.

% --- Project configuration ------------------------------------------------
\newcommand{\projectIssueBase}{https://github.com/LionSR/QICLean/issues/}
\newcommand{\projectPrBase}{https://github.com/LionSR/QICLean/pull/}
\newcommand{\projectDocBase}{https://sirui-lu.com/QICLean/docs/}
% ---------------------------------------------------------------------------

\DeclareUnicodeCharacter{2080}{\ifmmode{}_{0}\else\textsubscript{0}\fi}
\DeclareUnicodeCharacter{2081}{\ifmmode{}_{1}\else\textsubscript{1}\fi}
\DeclareUnicodeCharacter{2082}{\ifmmode{}_{2}\else\textsubscript{2}\fi}
\DeclareUnicodeCharacter{2099}{\ifmmode{}_{n}\else\textsubscript{n}\fi}

\newcommand{\Lean}{\textsc{Lean}}

% A formal declaration name, upright sans, breakable at underscores.
\ExplSyntaxOn
\NewDocumentCommand{\leanid}{m}{
  \ifmmode
    \textsf{\detokenize{#1}}
  \else
    \group_begin:
    \urlstyle{sf}
    \tl_set:Nn \l_tmpa_tl {#1}
    \tl_replace_all:Nnn \l_tmpa_tl {\_} {_}
    \exp_args:NV \path \l_tmpa_tl
    \group_end:
  \fi
}
\ExplSyntaxOff

% Traceability links; use in footnotes, not in the running prose.
\newcommand{\ghissue}[1]{\href{\projectIssueBase#1}{GitHub issue \##1}}
\newcommand{\ghpr}[1]{\href{\projectPrBase#1}{PR \##1}}
\newcommand{\doclink}[1]{\href{\projectDocBase#1.html}{\path{#1}}}

% Domain notation, using \providecommand so a note-local refinement stays
% possible.
\providecommand{\tr}{\operatorname{tr}}

% Dirac notation and the identity operator, as in blueprint/src/macros.
\usepackage{dsfont}
\providecommand{\ket}[1]{|#1\rangle}
\providecommand{\bra}[1]{\langle #1|}
\providecommand{\braket}[2]{\langle #1 | #2 \rangle}
\providecommand{\ketbra}[2]{|#1\rangle\!\langle #2|}
\providecommand{\Id}{\mathds{1}}
\providecommand{\one}{\mathds{1}}
\providecommand{\C}{\mathbb{C}}
\providecommand{\MN}[1]{M_{#1}(\C)}
\providecommand{\MD}{M_D(\C)}

% Verdict marker: \gapnote{<kind>}{<status>}. Kind is the policy
% classification (clarification, local-correction, scope-restriction,
% unfaithful, false-source, open-gap); status is open, wip (elimination
% underway), resolved, or historical. Machine-read by the site index;
% prints a verdict line.
\newcommand{\gapnote}[2]{\par\noindent\textbf{Verdict:} #1 (#2).\par}
. Project-wide notation
% belongs here, never in per-note redefinitions.

% --- Project configuration ------------------------------------------------
\newcommand{\projectIssueBase}{https://github.com/LionSR/QICLean/issues/}
\newcommand{\projectPrBase}{https://github.com/LionSR/QICLean/pull/}
\newcommand{\projectDocBase}{https://sirui-lu.com/QICLean/docs/}
% ---------------------------------------------------------------------------

\DeclareUnicodeCharacter{2080}{\ifmmode{}_{0}\else\textsubscript{0}\fi}
\DeclareUnicodeCharacter{2081}{\ifmmode{}_{1}\else\textsubscript{1}\fi}
\DeclareUnicodeCharacter{2082}{\ifmmode{}_{2}\else\textsubscript{2}\fi}
\DeclareUnicodeCharacter{2099}{\ifmmode{}_{n}\else\textsubscript{n}\fi}

\newcommand{\Lean}{\textsc{Lean}}

% A formal declaration name, upright sans, breakable at underscores.
\ExplSyntaxOn
\NewDocumentCommand{\leanid}{m}{
  \ifmmode
    \textsf{\detokenize{#1}}
  \else
    \group_begin:
    \urlstyle{sf}
    \tl_set:Nn \l_tmpa_tl {#1}
    \tl_replace_all:Nnn \l_tmpa_tl {\_} {_}
    \exp_args:NV \path \l_tmpa_tl
    \group_end:
  \fi
}
\ExplSyntaxOff

% Traceability links; use in footnotes, not in the running prose.
\newcommand{\ghissue}[1]{\href{\projectIssueBase#1}{GitHub issue \##1}}
\newcommand{\ghpr}[1]{\href{\projectPrBase#1}{PR \##1}}
\newcommand{\doclink}[1]{\href{\projectDocBase#1.html}{\path{#1}}}

% Domain notation, using \providecommand so a note-local refinement stays
% possible.
\providecommand{\tr}{\operatorname{tr}}

% Dirac notation and the identity operator, as in blueprint/src/macros.
\usepackage{dsfont}
\providecommand{\ket}[1]{|#1\rangle}
\providecommand{\bra}[1]{\langle #1|}
\providecommand{\braket}[2]{\langle #1 | #2 \rangle}
\providecommand{\ketbra}[2]{|#1\rangle\!\langle #2|}
\providecommand{\Id}{\mathds{1}}
\providecommand{\one}{\mathds{1}}
\providecommand{\C}{\mathbb{C}}
\providecommand{\MN}[1]{M_{#1}(\C)}
\providecommand{\MD}{M_D(\C)}

% Verdict marker: \gapnote{<kind>}{<status>}. Kind is the policy
% classification (clarification, local-correction, scope-restriction,
% unfaithful, false-source, open-gap); status is open, wip (elimination
% underway), resolved, or historical. Machine-read by the site index;
% prints a verdict line.
\newcommand{\gapnote}[2]{\par\noindent\textbf{Verdict:} #1 (#2).\par}
. Project-wide notation
% belongs here, never in per-note redefinitions.

% --- Project configuration ------------------------------------------------
\newcommand{\projectIssueBase}{https://github.com/LionSR/QICLean/issues/}
\newcommand{\projectPrBase}{https://github.com/LionSR/QICLean/pull/}
\newcommand{\projectDocBase}{https://sirui-lu.com/QICLean/docs/}
% ---------------------------------------------------------------------------

\DeclareUnicodeCharacter{2080}{\ifmmode{}_{0}\else\textsubscript{0}\fi}
\DeclareUnicodeCharacter{2081}{\ifmmode{}_{1}\else\textsubscript{1}\fi}
\DeclareUnicodeCharacter{2082}{\ifmmode{}_{2}\else\textsubscript{2}\fi}
\DeclareUnicodeCharacter{2099}{\ifmmode{}_{n}\else\textsubscript{n}\fi}

\newcommand{\Lean}{\textsc{Lean}}

% A formal declaration name, upright sans, breakable at underscores.
\ExplSyntaxOn
\NewDocumentCommand{\leanid}{m}{
  \ifmmode
    \textsf{\detokenize{#1}}
  \else
    \group_begin:
    \urlstyle{sf}
    \tl_set:Nn \l_tmpa_tl {#1}
    \tl_replace_all:Nnn \l_tmpa_tl {\_} {_}
    \exp_args:NV \path \l_tmpa_tl
    \group_end:
  \fi
}
\ExplSyntaxOff

% Traceability links; use in footnotes, not in the running prose.
\newcommand{\ghissue}[1]{\href{\projectIssueBase#1}{GitHub issue \##1}}
\newcommand{\ghpr}[1]{\href{\projectPrBase#1}{PR \##1}}
\newcommand{\doclink}[1]{\href{\projectDocBase#1.html}{\path{#1}}}

% Domain notation, using \providecommand so a note-local refinement stays
% possible.
\providecommand{\tr}{\operatorname{tr}}

% Dirac notation and the identity operator, as in blueprint/src/macros.
\usepackage{dsfont}
\providecommand{\ket}[1]{|#1\rangle}
\providecommand{\bra}[1]{\langle #1|}
\providecommand{\braket}[2]{\langle #1 | #2 \rangle}
\providecommand{\ketbra}[2]{|#1\rangle\!\langle #2|}
\providecommand{\Id}{\mathds{1}}
\providecommand{\one}{\mathds{1}}
\providecommand{\C}{\mathbb{C}}
\providecommand{\MN}[1]{M_{#1}(\C)}
\providecommand{\MD}{M_D(\C)}

% Verdict marker: \gapnote{<kind>}{<status>}. Kind is the policy
% classification (clarification, local-correction, scope-restriction,
% unfaithful, false-source, open-gap); status is open, wip (elimination
% underway), resolved, or historical. Machine-read by the site index;
% prints a verdict line.
\newcommand{\gapnote}[2]{\par\noindent\textbf{Verdict:} #1 (#2).\par}

\title{Relations among Relative Entropies: Support and Infinite Values}
\author{}
\date{2026-11-15}
\begin{document}
\maketitle
\gapnote{scope-restriction}{open}

\section{Source statement}
For positive semidefinite $\rho,\sigma\in M_d(\mathbb C)$ Wolf
\cite{Wolf2012Quantum}%
\footnote{\path{Notes/WolfNoteTexSource/ch08_distance_measures.tex},
lines~466--519, Equations~(8.33)--(8.39) and the footnote on line~474.}
sets
\begin{align}
    S^{(\alpha)}(\rho|\sigma)
        &= \frac{\tr[\rho-\rho^\alpha\sigma^{1-\alpha}]}{1-\alpha},
    &
    S(\rho|\sigma) &= \tr[\rho(\log\rho-\log\sigma)],
    \label{eq:rer_definitions}
\end{align}
and states that for all $\alpha\le1$ and $\beta\ge1$
\begin{align}
    S^{(\alpha)}(\rho|\sigma) &\le S(\rho|\sigma)\le S^{(\beta)}(\rho|\sigma).
    \label{eq:rer_source}
\end{align}
Outside $[0,1]$ the footnote takes inverses on the support and sets the
quantity to $+\infty$ when the kernel of the inverted operator is not
contained in the kernel of the other operator; $S(\rho|\sigma)=+\infty$
unless $\ker\sigma\subseteq\ker\rho$. At $\alpha=1$ and $\beta=1$ the
quantities are defined by their limits.

\section{Formal statement}
Powers are taken by spectral calculus with $x^\gamma=0$ at $x=0$ for
$\gamma\ne0$ and $X^0=\mathbb 1$; the logarithm is totalized by $\log0=0$.
Assume
\begin{align}
    \ker\sigma &\subseteq \ker\rho .
    \label{eq:rer_support}
\end{align}
With eigen-decompositions $\rho=\sum_i\lambda_i\lvert x_i\rangle\langle x_i\rvert$,
$\sigma=\sum_j\mu_j\lvert y_j\rangle\langle y_j\rvert$ and
$P_{ij}=\lvert\langle x_i|y_j\rangle\rvert^2$, every trace of the form
$\tr[f(\rho)g(\sigma)]$ equals $\sum_{ij}f(\lambda_i)g(\mu_j)P_{ij}$.
Condition (\ref{eq:rer_support}) gives $\lambda_iP_{ij}=0$ whenever
$\mu_j=0$. For $\lambda,\mu>0$ and every real $a$, the inequality
$\log x\le x-1$ at $x=(\mu/\lambda)^{1-a}$ gives
\begin{align}
    \lambda-\lambda^a\mu^{1-a}
        &\le (1-a)\,\lambda(\log\lambda-\log\mu),
    \label{eq:rer_scalar}
\end{align}
and the same inequality holds trivially when $\lambda=0$. Summing against
$P_{ij}$ yields
\begin{align}
    \tr[\rho-\rho^a\sigma^{1-a}] &\le (1-a)\,S(\rho|\sigma)
    \qquad (a\in\mathbb R).
    \label{eq:rer_common}
\end{align}
Division by $1-a$ gives (\ref{eq:rer_source}) for $\alpha<1<\beta$.

\section{Comparison with the source}
\begin{enumerate}
\item If (\ref{eq:rer_support}) fails, $S(\rho|\sigma)=+\infty$, and the
    right inequality reads $+\infty\le+\infty$ because $\sigma$ is then the
    inverted operator in $S^{(\beta)}$. The totalized trace-log expression is
    finite and does not represent $S$, so this case is excluded. It is a
    statement about extended real values that is not yet formalized.
\item The endpoints $\alpha=1$ and $\beta=1$ are excluded, since the quotient
    in (\ref{eq:rer_definitions}) is defined there only as a limit.
\item For $\alpha<0$ the powers of $\rho$ are taken with the
    support-inverse convention, without the footnote's assignment of
    $+\infty$. Under that assignment the left inequality is false as printed;
    this separate gap is recorded in
    \path{docs/paper-gaps/wolf_ch8_tsallis_negative_exponent_footnote.tex}.
\end{enumerate}
\doclink{QICLean/Analysis/RelativeEntropyRelations} proves
(\ref{eq:rer_common}) and both halves of (\ref{eq:rer_source}) under
(\ref{eq:rer_support}) for $\alpha<1<\beta$, without trace normalization.

\section{Verdict}
Scope restriction to finite relative entropy and to $\alpha\ne1\ne\beta$.
Elimination: state (\ref{eq:rer_source}) with extended-real values,
including the infinite cases and the limit endpoints; tracked in
\ghissue{46}.

\bibliographystyle{plain}
\bibliography{references}
\end{document}
